\chapter{Experiment Tracking and Reproducibility}\label{ch:ml:experiment-tracking-and-reproducibility}

A regulator asks a firm to show the model that traded on a given morning, the data it was trained on and every alternative
it was chosen over. The firm can name the model; it takes three weeks to retrain something close to it, and it cannot say how
many alternatives there were. This chapter builds the records that make the answer a query: a tracker that stores every run
with everything that determined it, a registry that says which model was in production when and where it came from, and an
\omterm{def:ml:experiment-tracking-and-reproducibility:registry}{audit trail} that cannot be edited without the edit showing. On the chapter's study, 200 tuning runs of a boosted forecast, the
registry reproduces the March model bit for bit from its records alone; with any one of four fields missing, nothing close to
it comes back. The tracker also answers the question the regulator did not ask: counting the 200 trials, the champion's
validation Sharpe ratio of 2.22 has a deflated Sharpe ratio of 0.15.

\section{What must be recorded}

\begin{definition}[Experiment tracker, run record]\label{def:ml:experiment-tracking-and-reproducibility:tracker}
An \emph{experiment tracker}\index{experiment tracker} stores a record of every training run a team makes, including the
ones that are discarded. A \emph{run record}\index{run record} holds
everything that determines the run's result (parameters, seeds, the data snapshot, the code version, the environment) and
everything the run produced (metrics, artefacts, logs), with the time and the author.
\end{definition}

The research log of Book~7 (chapter~1) recorded trials to count them; the tracker records them to rebuild them. The chapter's
tracker (\cref{lst:ml:exp:run}) writes each run as a JSON record whose identifier is the hash of its content, stores the fitted
model once under the content hash of its dump (Book~7's \texttt{firm.workflow}), and logs the run as a trial in
\texttt{firm.researchlog}. Timestamps are passed in, never read from the clock, so that the tracker's own contents are
reproducible.

The study: a gradient-boosted forecast (chapter~5) of next-day returns on a synthetic panel of 40 names and eight features over
700 days, trained on the first 400, scored by the Sharpe ratio of a daily long-short portfolio on the next 150, and tested on the
last 150. Two hundred runs draw five parameters and a seed at random. Their validation Sharpe ratios range from $-3.16$ to 2.22,
with a median of $-0.83$ (\cref{fig:ml:exp:sharpes}): the signal is weak, and most configurations fit noise.

\begin{omfigure}
\begin{tikzpicture}
\begin{axis}[width=11cm,height=5.8cm,ybar,bar width=12pt,xlabel={validation Sharpe ratio (annualised)},ylabel={runs},
  tick label style={font=\small},label style={font=\small},xmin=-4,xmax=4,ymin=0,ymax=42,grid=major,grid style={gray!25},
  table/col sep=comma]
\addplot[fill=omProp!60,draw=none] table[x=x,y=runs]{figdata/ml/25-experiment-tracking-and-reproducibility/sharpes.csv};
\addplot[ycomb,omThm,very thick,mark=*,restrict expr to domain={\coordindex}{0:0}] table[x=x,y expr=40]{figdata/ml/25-experiment-tracking-and-reproducibility/marks.csv};
\addplot[ycomb,black,very thick,dashed,mark=none,restrict expr to domain={\coordindex}{1:1}] table[x=x,y expr=40]{figdata/ml/25-experiment-tracking-and-reproducibility/marks.csv};
\addplot[ycomb,omDef,very thick,mark=square*,restrict expr to domain={\coordindex}{2:2}] table[x=x,y expr=40]{figdata/ml/25-experiment-tracking-and-reproducibility/marks.csv};
\end{axis}
\end{tikzpicture}

\omcaption{Validation Sharpe ratios of the 200 tracked runs; circle: the champion's (2.22); square: its test Sharpe ratio
(0.91); dashed: the expected maximum of 200 independent strategies with no skill over 150 days (3.58). Data:
\texttt{ml\_track.study}.}\label{fig:ml:exp:sharpes}
\end{omfigure}

The champion's test Sharpe ratio is 0.91, less than half its validation figure. Its probabilistic Sharpe ratio, the
probability that its true Sharpe ratio is positive as if it were the only run, is 0.96; its deflated Sharpe ratio (Book~4,
chapter~12), with the trial count taken from the tracker, is 0.15. The expected best of 200 skill-less strategies over 150 days
is 3.58, above the champion: the correction here is conservative, because the runs are correlated and are fewer independent
trials than 200, but a result that cannot beat the best of its own noise has not earned production on validation evidence
alone.

\section{Versioning data, code and models}

\begin{definition}[Data versioning, model artefact]\label{def:ml:experiment-tracking-and-reproducibility:versions}
\emph{Data versioning}\index{data versioning} identifies every dataset a run reads by an immutable snapshot (a content hash or a
dated, frozen version), so that later corrections create new versions instead of changing old ones. A \emph{model
artefact}\index{model artefact} is the stored result of a training run, the fitted parameters in a portable form, identified by
its content hash.
\end{definition}

Versions exist because data and code change under a model. In the study, the data vendor revises 1\% of the returns in March
(snapshot 2026-03 replaces 2026-01), and the feature code changes its clipping from three standard deviations to two and a half
(v2 replaces v1). The champion was trained on 2026-01 with v1. The record says so; without it, a team would use what it has now.

\section{The registry and the audit trail}

\begin{definition}[Model registry, model lineage, audit trail]\label{def:ml:experiment-tracking-and-reproducibility:registry}
A \emph{model registry}\index{model registry} lists the firm's models by name and version with their stage (candidate,
production, retired) and the history of stage changes. \emph{Model lineage}\index{model lineage} links a registered version to
the run, data snapshot and code that produced it. An \emph{audit trail}\index{audit trail} is an append-only record of actions
(registrations, promotions, overrides) with who, when and why, protected against silent edits; here each entry carries the hash
of the previous one (\cref{lst:ml:exp:audit}).
\end{definition}

\begin{dated}{2026-09}{Records of algorithmic trading changes in the EU}\label{dat:ml:experiment-tracking-and-reproducibility:rts6}
Commission Delegated Regulation (EU) 2017/589 (the regulatory technical standards on organisational requirements for
algorithmic trading, known as RTS~6) requires an investment firm to ``keep records of any material change made to the software
used for algorithmic trading'', allowing it to determine when a change was made, who made it, who approved it and its nature
(Article 5(7)); order records of high-frequency traders are kept for five years (Article 28).
\end{dated}

The registry answers the regulator's first question by time (\cref{lst:ml:exp:asof}): the version in production on the morning
of 10 March is found from the stage histories, and its lineage leads to the \omterm{def:ml:experiment-tracking-and-reproducibility:tracker}{run record}. The \omterm{def:ml:experiment-tracking-and-reproducibility:registry}{audit trail} answers ``who promoted it,
and why'': version 1 of \texttt{xs-gbdt}, promoted on 2 March as the champion of 200 runs. Changing any character of any line of
the trail breaks the hash chain at that line.

\section{Reproducing a model, and when bitwise is too much}

From the registry entry alone the chapter re-runs the training (\cref{lst:ml:exp:repro}) and obtains a model whose dump has the
same content hash as the stored artefact: bit for bit the model that traded. Then it removes one field of the record at a time
and replaces it with what a team would assume without it (\cref{tab:ml:exp:ablation}). Every removal breaks reproduction, and
none is harmless: the reproduced models' test Sharpe ratios differ from the original's by $-1.28$ with default parameters,
$-1.25$ with the revised data, $-0.44$ with the current feature code, and $+0.85$ with the default seed.

\begin{table}[htbp]
\centering\small
\begin{tabular}{@{}llcc@{}}
\toprule
field missing & assumed instead & reproduced bit for bit & change in test Sharpe ratio\\
\midrule
none (the registry's record) & & yes & 0\\
parameters & the firm's defaults & no & $-1.28$\\
seed & 1 & no & $+0.85$\\
data snapshot & the latest (2026-03) & no & $-1.25$\\
code version & the current (v2) & no & $-0.44$\\
\bottomrule
\end{tabular}
\caption{Reproducing the production model from its record, with each field in turn replaced by an assumption. Data:
\texttt{ml\_track.reproduction}.}\label{tab:ml:exp:ablation}
\end{table}

The seed's row carries a second lesson: a different seed of the same configuration moves the test Sharpe ratio by 0.85, as
much as the difference between configurations. A result that does not survive a change of seed was not a result, and reporting
the spread over seeds (chapter~7) is part of what a tracker enables.

Reproduction bit for bit is the right standard for a model's own artefact, rebuilt on the same software and hardware; the chapter's
fits are deterministic because LightGBM runs single-threaded with its deterministic option. Across library versions, thread
counts or hardware (chapter~23's \omterm{def:ml:training-infrastructure:mixed}{bfloat16}), floating-point sums change order and bits change; there the standard is statistical:
the same metrics within the spread over seeds, checked by re-running. What must always be exact is the record: the data, code,
parameters and seed that produced the artefact.

\begin{method}[Tracking research that trades]\label{met:ml:exp:method}
\begin{enumerate}
\item Track every run, discarded ones included, with its parameters, seed, data snapshot, code version and environment, and
store artefacts by content hash.
\item Take the trial count for deflation from the tracker, not from memory.
\item Register production models with lineage and stage history; promote and retire only through the registry, with reasons in
an append-only trail.
\item Test reproduction from the registry alone, regularly, and treat a failure as an incident.
\end{enumerate}
\end{method}

\section{Tutorial: reproduce the March model}\label{tut:ml:experiment-tracking-and-reproducibility:tut}

\begin{tutorial}
\textbf{Goal.} Track 200 tuning runs, register and promote the champion, find it again by date, reproduce it bitwise, and break
reproduction one field at a time. \textbf{End state:} \cref{fig:ml:exp:sharpes}, \cref{tab:ml:exp:ablation}.
\begin{tutsteps}
\item \textbf{A \omterm{def:ml:experiment-tracking-and-reproducibility:tracker}{run record}.}
\omcode{code/firm/exptrack/firm_exptrack.py}{73}{85}{Recording a run: artefact by content hash, record by its own hash, trial in the research log.}\label{lst:ml:exp:run}
\item \textbf{An append-only trail.}
\omcode{code/firm/exptrack/firm_exptrack.py}{48}{64}{Hash-chained audit entries and their verification.}\label{lst:ml:exp:audit}
\item \textbf{The model in production on a given date.}
\omcode{code/firm/exptrack/firm_exptrack.py}{134}{144}{The registry queried as of a time.}\label{lst:ml:exp:asof}
\item \textbf{Reproduction.}
\omcode{code/firm/exptrack/firm_exptrack.py}{151}{155}{Re-running from a record and comparing hashes.}\label{lst:ml:exp:repro}
\item \textbf{Run} \texttt{ml\_track.summary()}, \texttt{reproduction()}, \texttt{null\_max()} and \texttt{fig\_track.py}
(about 45 seconds on one core).
\end{tutsteps}
\textbf{What to change next.} Estimate the effective number of independent trials from the correlation of the runs' validation
returns, and recompute the deflated Sharpe ratio with it.
\end{tutorial}

\section{Build: experiment tracking}\label{bld:ml:experiment-tracking-and-reproducibility:exptrack}

\begin{build}
\bfield{Purpose} Every model the firm trades traceable to its run, data and code, and rebuildable from them.
\bfield{Interface} \texttt{Tracker(root, log, family)} with \texttt{run}, \texttt{get}, \texttt{artefact}, \texttt{runs},
\texttt{best}; \texttt{Registry(root, audit)} with \texttt{register}, \texttt{promote}, \texttt{current}, \texttt{as\_of},
\texttt{lineage}; \texttt{AuditTrail(path)} with \texttt{append} and \texttt{verify}; \texttt{reproduce(record, train)}; on
\texttt{firm.workflow} and \texttt{firm.researchlog}.
\bfield{Rules} Records are written once; artefacts are addressed by content; the clock never enters a record except as a
supplied timestamp; promotions carry reasons.
\bfield{Acceptance tests} \texttt{code/firm/exptrack/tests/}: identical runs share an artefact; the registry returns the version
in production at a past time; promoting a new version retires the old; a tampered audit line is found; reproduction succeeds
from a complete record and fails when a field changes.
\bfield{Stretch} Environment capture (package versions, hardware) and a check that refuses to reproduce across a different
environment; data snapshots on \texttt{firm.pit}.
\end{build}

\begin{omsources}
\item J.~Pineau and co-authors, ``Improving reproducibility in machine learning research (a report from the NeurIPS 2019
reproducibility program)'', arXiv:2003.12206, 2020.
\item O.~E.~Gundersen and S.~Kjensmo, ``State of the art: reproducibility in artificial intelligence'', \emph{AAAI}, 2018.
\item Commission Delegated Regulation (EU) 2017/589 of 19 July 2016 (RTS~6), \emph{Official Journal of the European Union}.
\end{omsources}

\section{Exercises}

\begin{exercise}[$\star$]\label{exo:ml:experiment-tracking-and-reproducibility:1}
List the five fields of a \omterm{def:ml:experiment-tracking-and-reproducibility:tracker}{run record} that determine its result, and one thing each can silently change.
\end{exercise}

\begin{exercise}[$\star$]\label{exo:ml:experiment-tracking-and-reproducibility:2}
Why is a run's identifier the hash of its record rather than a counter or a timestamp?
\end{exercise}

\begin{exercise}[$\star$]\label{exo:ml:experiment-tracking-and-reproducibility:3}
What does the deflated Sharpe ratio of 0.15 mean, and why does it differ from the probabilistic Sharpe ratio of 0.96?
\end{exercise}

\begin{exercise}[$\star\star$]\label{exo:ml:experiment-tracking-and-reproducibility:4}
Why is the expected maximum of 200 null Sharpe ratios too harsh a benchmark for these 200 runs?
\end{exercise}

\begin{exercise}[$\star\star$]\label{exo:ml:experiment-tracking-and-reproducibility:5}
How does a hash chain make an \omterm{def:ml:experiment-tracking-and-reproducibility:registry}{audit trail} tamper-evident, and what does it not protect against?
\end{exercise}

\begin{exercise}[$\star\star$]\label{exo:ml:experiment-tracking-and-reproducibility:6}
\emph{Find the flaw.} ``We save every production model's weights, so we can always reproduce our models.''
\end{exercise}

\begin{exercise}[$\star\star\star$]\label{exo:ml:experiment-tracking-and-reproducibility:7}
\emph{Coding.} Retrain the champion's configuration with ten different seeds and report the spread of test Sharpe ratios.
What does it say about the champion?
\end{exercise}

\begin{exercise}[$\star\star\star$]\label{exo:ml:experiment-tracking-and-reproducibility:8}
Design what the registry and the \omterm{def:ml:experiment-tracking-and-reproducibility:registry}{audit trail} must contain for a firm to satisfy RTS~6's record of material changes for a model
that is retrained weekly.
\end{exercise}

\section{Problem: Reproduce the March Model}

\begin{problem}[{Weekend problem --- the model that traded}]\label{pb:ml:experiment-tracking-and-reproducibility:1}
The chapter's 200 runs, registry and trail.

\medskip\noindent\textbf{Part I --- The study.}
\begin{enumerate}
\item What does each \omterm{def:ml:experiment-tracking-and-reproducibility:tracker}{run record} contain?
\item What are the distribution of validation Sharpe ratios and the champion's figures?
\item What are its probabilistic and deflated Sharpe ratios?
\item What is the expected best of 200 null strategies, and why is it above the champion?
\end{enumerate}

\medskip\noindent\textbf{Part II --- Versions and registry.}
\begin{enumerate}[resume]
\item What changed in the data and the code after the champion was trained?
\item How is the model in production on 10 March found?
\item What does the \omterm{def:ml:experiment-tracking-and-reproducibility:registry}{audit trail} record, and how is tampering detected?
\item What does RTS~6 require the firm to record?
\end{enumerate}

\medskip\noindent\textbf{Part III --- Reproduction.}
\begin{enumerate}[resume]
\item What happens when the model is rebuilt from its record?
\item What happens when each field is missing?
\item What does the seed's row show?
\item When is bitwise reproduction the wrong standard?
\end{enumerate}

\medskip\noindent\textbf{Part IV --- The verdict.}
\begin{enumerate}[resume]
\item State the \emph{named result}: the number of runs behind the champion and its deflated Sharpe ratio, and the field whose
absence made reproduction fail.
\item Would you have promoted the champion? What would you have asked for?
\item How would you answer the regulator's request with these tools, and how long would it take?
\item What else would the tracker need to satisfy a five-year retention rule?
\item How should the tracker count trials that were run but never recorded?
\item What should be in the registry for a model retrained automatically every week?
\item Who should be allowed to promote a model?
\item In one sentence: what does it take to reproduce a model?
\end{enumerate}
\end{problem}

\section{Interview questions}

\begin{interviewq}[$\star$ \iqroles{mle}]\label{iq:ml:experiment-tracking-and-reproducibility:1}
What does an \omterm{def:ml:experiment-tracking-and-reproducibility:tracker}{experiment tracker} record, and why record the failed runs?
\end{interviewq}

\begin{interviewq}[$\star\star$ \iqroles{mle, developer}]\label{iq:ml:experiment-tracking-and-reproducibility:2}
How do you make a training run reproducible bit for bit?
\end{interviewq}

\begin{interviewq}[$\star\star$ \iqroles{researcher}]\label{iq:ml:experiment-tracking-and-reproducibility:3}
How does the number of trials change what a backtest's Sharpe ratio means?
\end{interviewq}

\begin{interviewq}[$\star\star$ \iqroles{mle}]\label{iq:ml:experiment-tracking-and-reproducibility:4}
Design a \omterm{def:ml:experiment-tracking-and-reproducibility:registry}{model registry} for a trading firm. What stages and fields does it need?
\end{interviewq}

\begin{interviewq}[$\star\star$ \iqroles{mle, developer}]\label{iq:ml:experiment-tracking-and-reproducibility:5}
Why version data, and how would you version a dataset that is corrected after the fact?
\end{interviewq}

\begin{interviewq}[$\star\star\star$ \iqroles{mle}]\label{iq:ml:experiment-tracking-and-reproducibility:6}
Your production model cannot be reproduced. Walk through the investigation.
\end{interviewq}
